% Generated draft: requires mathematical and source review.
\documentclass[11pt]{article}
\usepackage[margin=25mm]{geometry}
\usepackage{fontspec}
\setmainfont{Latin Modern Roman}
\usepackage{amsmath,amssymb,mathtools,unicode-math}
\setmathfont{Latin Modern Math}
\usepackage{enumitem,needspace,xurl,hyperref}
\setlength{\parindent}{0pt}
\setlength{\parskip}{5pt}
\setlength{\emergencystretch}{3em}
\newcommand{\N}{\mathbb{N}}
\newcommand{\Z}{\mathbb{Z}}
\newcommand{\Q}{\mathbb{Q}}
\newcommand{\R}{\mathbb{R}}
\newcommand{\C}{\mathbb{C}}
\newcommand{\F}{\mathbb{F}}
\newcommand{\A}{\mathbb{A}}
\newcommand{\PP}{\mathbb{P}}
\newcommand{\E}{\mathbb{E}}
\newcommand{\eps}{\varepsilon}
\newcommand{\dd}{\,\mathrm{d}}
\begin{document}
\section*{UnsolvedMath ID: 11000083}
\textbf{Source problem code:} catalogue entry
\subsection{Definitions and mathematical statement}
Let \(\mathcal M_g\) be the moduli orbifold with its Teichm\"uller metric. Let \(\mathscr G_{\mathrm{simple}}(\mathcal M_g)\) denote the class of closed geodesics meant by ``simple length spectrum'' in the source, with its intended equivalence and multiplicity convention. For \(t>0\), put
\[
 m_g(t)=\#\{[c]\in\mathscr G_{\mathrm{simple}}(\mathcal M_g):
 \ell_{\mathrm{Teich}}(c)=t\}.
\]
The excerpt does not define ``simple,'' whether iterates or orientations are distinguished, or the equivalence on geodesics. Geometric simplicity and primitivity are different restrictions and must not be interchanged without the source definition.

\emph{Question 7.7 (simple length spectrum).} Does the simple length spectrum of \(\mathcal M_g\) have bounded multiplicity? In the notation above, does there exist a finite \(B(g)\), independent of \(t\), such that
\[
 m_g(t)\leq B(g)\qquad\text{for every }t>0?
\]
If so, determine how such a bound depends on genus. The question is about the Teichm\"uller metric on moduli space, not merely the length spectrum of simple curves on one fixed hyperbolic surface. The source states that it includes the corresponding bounded-multiplicity question for many hyperbolic surfaces and describes that related question as open in its context, but supplies no construction or reduction. Its exact intended geodesic notion and genus range must be recovered before this becomes a fixed counting problem.
\subsection{Short English statement}
For moduli space with the Teichm\"uller metric, ask whether each simple-geodesic length occurs at most a genus-dependent number of times, uniformly over lengths. The meaning of simple and the counting convention are omitted and must not be replaced by a different surface-spectrum question.
\subsection{Sources}
\begin{itemize}
\item Local UnsolvedMath registry, permanent ID 11000083, source code catalogue entry, Question 7.7, ``Simple length spectrum,'' PDF page 58. No individual proposer is supplied.
\item Farb, \emph{Problems on Mapping Class Groups and Related Topics}. Supplied URL: \url{https://www.math.uchicago.edu/~farb/papers/mcgbook.pdf}; recorded access 2026-07-29; extraction pdf-font-structure. The edited volume contains 329 explicitly typeset headings across chapters, with repeated local numbering disambiguated by page.
\item The source makes a connection to the corresponding question for many hyperbolic surfaces without a supplied citation or reduction. The research summary mentions Wolpert and classical simple closed-geodesic spectra, while explicitly saying the exact conjecture is unclear; no full reference is given.
\item Supplied status open and OPEN-TRIAGE. Historical related-question openness and the summary discussion are not independently verified current facts. Current primary-literature status and redistribution terms for the publicly accessible author/publisher copy remain NEEDS\_REVIEW.
\end{itemize}
\end{document}
